\chapter{Exception Handling in Python: Errors and Safe Programming}

\section{Introduction: Real-Life Hook}

Students encounter mistakes every day:  

\begin{itemize}
    \item Forgetting homework or submitting the wrong file  
    \item Dividing by zero in math homework  
    \item Searching for a contact that does not exist  
\end{itemize}

Similarly, \textbf{programs can have errors}. Python provides \textbf{exception handling} to \textbf{catch errors and prevent programs from crashing}.

\section{Concept: Errors and Exceptions}

\begin{definition}{Exception}
An \textbf{exception} is an error detected during program execution that stops the program unless handled.
\end{definition}

\subsection{Example: Division by Zero}

\begin{verbatim}
num = 5
den = 0
print(num / den)  # This raises a ZeroDivisionError
\end{verbatim}

- This will crash the program if not handled  

\section{Try and Except Block}

\begin{verbatim}
try:
    num = 5
    den = 0
    print(num / den)
except ZeroDivisionError:
    print("Cannot divide by zero!")
\end{verbatim}

- \texttt{try} contains code that might raise an error  
- \texttt{except} handles the error gracefully  

\subsection{Catching Multiple Exceptions}

\begin{verbatim}
try:
    x = int(input("Enter a number: "))
    y = int(input("Enter another number: "))
    print(x / y)
except ZeroDivisionError:
    print("Cannot divide by zero!")
except ValueError:
    print("Invalid input, enter numbers only!")
\end{verbatim}

- Multiple \texttt{except} blocks handle different types of errors  

\section{Using Else and Finally}

\begin{verbatim}
try:
    num = int(input("Enter a number: "))
    print("Number squared:", num ** 2)
except ValueError:
    print("Invalid input!")
else:
    print("Input was valid.")
finally:
    print("This runs no matter what.")
\end{verbatim}

- \texttt{else} runs if no exception occurs  
- \texttt{finally} always runs (e.g., closing a file or cleaning up)  

\section{Mini-Project: Safe Division Program}

Create a program that takes two numbers from the user and divides them safely:

\begin{verbatim}
try:
    x = float(input("Enter numerator: "))
    y = float(input("Enter denominator: "))
    result = x / y
except ZeroDivisionError:
    print("Cannot divide by zero!")
except ValueError:
    print("Invalid input! Enter numbers only.")
else:
    print("Result:", result)
finally:
    print("Program ended safely.")
\end{verbatim}

\section{Step-by-Step Diagram: Exception Handling Flow}

\begin{center}
\begin{tikzpicture}[node distance=1.5cm]
\node (start) [draw, ellipse, fill=green!20] {Start Program};
\node (try) [draw, rectangle, fill=yellow!20, below of=start] {Execute Code in Try Block};
\node (exception) [draw, diamond, fill=orange!20, below of=try] {Exception Occurred?};
\node (except) [draw, rectangle, fill=yellow!20, below left of=exception, xshift=-2cm] {Handle Exception in Except Block};
\node (else) [draw, rectangle, fill=yellow!20, below right of=exception, xshift=2cm] {Run Else Block if No Exception};
\node (finally) [draw, rectangle, fill=orange!20, below of=exception, yshift=-3cm] {Finally Block Always Executes};
\node (end) [draw, ellipse, fill=green!20, below of=finally] {End Program};
\draw[->] (start) -- (try);
\draw[->] (try) -- (exception);
\draw[->] (exception.west) -| node[near start, above] {Yes} (except.north);
\draw[->] (exception.east) -| node[near start, above] {No} (else.north);
\draw[->] (except.south) |- (0,-6.5) -| (finally.north);
\draw[->] (else.south) |- (0,-6.5) -| (finally.north);
\draw[->] (finally.south) -- (end.north);
\end{tikzpicture}
\end{center}

\section{Key Concepts}

\begin{itemize}
    \item Exceptions are runtime errors that can stop programs  
    \item \texttt{try} block contains risky code  
    \item \texttt{except} block handles errors gracefully  
    \item \texttt{else} runs if no exception occurs; \texttt{finally} always runs  
    \item Common exceptions: \texttt{ZeroDivisionError}, \texttt{ValueError}, \texttt{FileNotFoundError}  
\end{itemize}

\section{End-of-Chapter Exercises}

\begin{exercise}{Very Short Questions}
\begin{enumerate}
    \item What is an exception in Python?  
    \item Name two common exceptions.  
    \item What is the purpose of a \texttt{finally} block?  
\end{enumerate}
\end{exercise}

\begin{exercise}{Short Answer Questions}
\begin{enumerate}
    \item Write code to safely divide two numbers entered by the user.  
    \item Explain the difference between \texttt{except} and \texttt{finally}.  
    \item How can you catch multiple exceptions in a program?  
\end{enumerate}
\end{exercise}

\begin{exercise}{Activity-Based Questions}
\begin{enumerate}
    \item Write a program to open a file and read its content safely. Handle \texttt{FileNotFoundError}.  
    \item Create a program that asks the user for an integer and prints its square. Handle \texttt{ValueError} for invalid input.  
    \item Make a mini-calculator that can add, subtract, multiply, or divide two numbers safely with exception handling.  
\end{enumerate}
\end{exercise}

\section{Teacher Notes}

\marginnote{
\textbf{Learning Objectives:}  
\begin{itemize}
    \item Students understand what exceptions are and how to handle them  
    \item Students can use \texttt{try}, \texttt{except}, \texttt{else}, and \texttt{finally} to write safe programs  
    \item Students can handle common runtime errors gracefully  
\end{itemize}

\textbf{Teaching Tips:}  
\begin{itemize}
    \item Demonstrate each exception with a real-life analogy (e.g., dividing by zero)  
    \item Encourage experimenting with multiple types of exceptions  
    \item Show how \texttt{finally} can be used to clean up resources (like closing files)  
\end{itemize}

\textbf{Assessment:}  
\begin{itemize}
    \item Check correct usage of try-except-else-finally blocks  
    \item Observe understanding of exception types and safe programming practices  
\end{itemize}
}